\section{Introduction}
\label{sec:intro}

Serverless platforms run small pieces of customer code, called \emph{functions}, whenever a request arrives. Most of these functions finish in well under a second~\cite{shahrad2020wild}, so the useful work is tiny compared with everything around it.
Before a function can run, the platform must create an isolated environment for it and load its language runtime. This is the \emph{cold start}. The request must also travel to wherever the function happens to run.
Providers hide the cold start by keeping environments alive in their datacenters, which holds memory around the clock~\cite{shahrad2020wild,fuerst2021faascache}. They can do little about distance.

Meanwhile, millions of home and office computers are close to users and idle for most of the day. Volunteer computing has used such machines for two decades~\cite{anderson2002seti,anderson2004boinc}, but only for batch jobs with deadlines of hours.
We ask a harder question: can volunteer machines answer requests that must finish in a few seconds?

Two things stand in the way.
The first is \textbf{churn}. An owner can switch off, reboot or reclaim the machine at any moment, and everything the platform prepared on it disappears: warm environments, cached code and the secret keys it was given. For a function that runs for tens of milliseconds, preparing this state again costs far more than the function.
The second is \textbf{distrust}. The owner controls the hardware and the operating system, so the owner can read the customer's code and data, return forged results, or walk away at the worst moment.
The two feed on each other. The usual cure for distrust is to check a machine before using it, and that check takes time. Churn throws the result away each time a session ends.

Earlier work solves only parts of this. Volunteer-computing systems tolerate untrusted machines by running every task several times and comparing answers~\cite{anderson2004boinc}. That multiplies cost, works only for tasks with a single right answer, and still shows the input to whoever runs it.
Serverless schedulers avoid cold starts well~\cite{shahrad2020wild,fuerst2021faascache} but assume trusted machines that stay up. Confidential serverless designs protect customers from the platform operator~\cite{arnautov2016scone,wallet2025} but run on provider servers and treat the trust check as a one-time setup.
And the standard way to notice a dead machine, listening for its periodic ``I am alive'' message (a heartbeat), fails against an owner who keeps sending heartbeats while throwing work away. Such a quiet failure is a \emph{gray failure}~\cite{huang2017gray}.

\textbf{Our starting point} is that proving what software a machine runs changes the fault model. If the customer's code and the scheduler's own code run only inside a locked, measured software image, and a verifier has checked hardware-signed evidence of that image (\emph{remote attestation}), the owner cannot read, forge or tamper with anything inside.
The owner keeps only the power switch and the network cable. A hostile owner can therefore do exactly what churn already does: stop, or delay. This failure class has a standard name, \emph{omission and timing failures}~\cite{cristian1991understanding,avizienis2004basic}, and a volunteer platform has to tolerate it anyway. One set of mechanisms then covers crashes, churn and malice.

That reduction exposes the real scheduling problem. A machine is \emph{attested} for one session. It becomes \emph{provisioned} for a customer once that customer's key service has released a secret key to it, and \emph{warm} for a function once a ready environment is waiting.
We call these together \emph{attested session state}. It is \emph{slow to create}: on our measurements a powered-off machine needs \CalAttested\,s to become an attested worker, thousands of times longer than resuming a warm environment. It is \emph{perishable}, since every session end destroys it. And it is \emph{exposing}, since each machine holding a customer's key is one more place the key could leak.
Choosing where this state lives is a scheduling decision with a price, and it cannot be left to the per-request path.

\voss{} is a scheduler built around this view (Figure~\ref{fig:arch}). It has three ideas. First, it separates \emph{where state lives} from \emph{where a request runs}. A slow loop decides which machines hold which keys and warm environments, with spare capacity so that losing a few machines does not hurt. A fast loop then picks a machine for each request and duplicates the risky ones.
Second, it detects failures from missing \emph{progress}, meaning work accepted but never finished, instead of missing heartbeats, and it replaces a failed scheduler without asking volunteers to agree on anything. Third, it records work with signed receipts rather than a blockchain.
Because volunteer machines sit behind home routers and change address often, machines reach the scheduler through a provider-hosted publish/subscribe broker with one topic per cluster (Section~\ref{sec:network}).

Our contributions are:
(1)~showing that attestation reduces hostile owners to omission and timing failures, and stating the resulting problem of managing attested session state;
(2)~a two-timescale scheduler whose planner keeps coverage even if any $f$ machines fail together;
(3)~a communication architecture for unstable, unreachable volunteer addresses, and a comparison of volunteer-hosted and provider-hosted schedulers;
(4)~fault tolerance without agreement among volunteers;
(5)~a deterministic, calibrated evaluation with hypotheses and decision rules fixed in advance, including what did not work.

Planning is the largest lever: it lowers tail latency by \GreedyNoneMed--\GreedyNoneHigh\% against no planning, and progress-based detection keeps \VossSuccAlpha\% of requests succeeding when 20\% of owners drop work, where heartbeat-only detection falls to \HbSuccAlpha\%.
\voss{} meets a 3\,s tail-latency target in most settings we map, and the strongest simple baseline never does. Yet some designed components do not earn their complexity: geographic clustering was slower than one scheduler, an exact solver was not distinguishable from a greedy heuristic, and a commercial-style reference is still faster. All results come from simulation and one calibration machine (Section~\ref{sec:threats}).
